% Options for packages loaded elsewhere
\PassOptionsToPackage{unicode}{hyperref}
\PassOptionsToPackage{hyphens}{url}
\documentclass[
]{article}
\usepackage{xcolor}
\usepackage[margin=1in]{geometry}
\usepackage{amsmath,amssymb}
\setcounter{secnumdepth}{-\maxdimen} % remove section numbering
\usepackage{iftex}
\ifPDFTeX
  \usepackage[T1]{fontenc}
  \usepackage[utf8]{inputenc}
  \usepackage{textcomp} % provide euro and other symbols
\else % if luatex or xetex
  \usepackage{unicode-math} % this also loads fontspec
  \defaultfontfeatures{Scale=MatchLowercase}
  \defaultfontfeatures[\rmfamily]{Ligatures=TeX,Scale=1}
\fi
\usepackage{lmodern}
\ifPDFTeX\else
  % xetex/luatex font selection
\fi
% Use upquote if available, for straight quotes in verbatim environments
\IfFileExists{upquote.sty}{\usepackage{upquote}}{}
\IfFileExists{microtype.sty}{% use microtype if available
  \usepackage[]{microtype}
  \UseMicrotypeSet[protrusion]{basicmath} % disable protrusion for tt fonts
}{}
\makeatletter
\@ifundefined{KOMAClassName}{% if non-KOMA class
  \IfFileExists{parskip.sty}{%
    \usepackage{parskip}
  }{% else
    \setlength{\parindent}{0pt}
    \setlength{\parskip}{6pt plus 2pt minus 1pt}}
}{% if KOMA class
  \KOMAoptions{parskip=half}}
\makeatother
\usepackage{color}
\usepackage{fancyvrb}
\newcommand{\VerbBar}{|}
\newcommand{\VERB}{\Verb[commandchars=\\\{\}]}
\DefineVerbatimEnvironment{Highlighting}{Verbatim}{commandchars=\\\{\}}
% Add ',fontsize=\small' for more characters per line
\usepackage{framed}
\definecolor{shadecolor}{RGB}{248,248,248}
\newenvironment{Shaded}{\begin{snugshade}}{\end{snugshade}}
\newcommand{\AlertTok}[1]{\textcolor[rgb]{0.94,0.16,0.16}{#1}}
\newcommand{\AnnotationTok}[1]{\textcolor[rgb]{0.56,0.35,0.01}{\textbf{\textit{#1}}}}
\newcommand{\AttributeTok}[1]{\textcolor[rgb]{0.13,0.29,0.53}{#1}}
\newcommand{\BaseNTok}[1]{\textcolor[rgb]{0.00,0.00,0.81}{#1}}
\newcommand{\BuiltInTok}[1]{#1}
\newcommand{\CharTok}[1]{\textcolor[rgb]{0.31,0.60,0.02}{#1}}
\newcommand{\CommentTok}[1]{\textcolor[rgb]{0.56,0.35,0.01}{\textit{#1}}}
\newcommand{\CommentVarTok}[1]{\textcolor[rgb]{0.56,0.35,0.01}{\textbf{\textit{#1}}}}
\newcommand{\ConstantTok}[1]{\textcolor[rgb]{0.56,0.35,0.01}{#1}}
\newcommand{\ControlFlowTok}[1]{\textcolor[rgb]{0.13,0.29,0.53}{\textbf{#1}}}
\newcommand{\DataTypeTok}[1]{\textcolor[rgb]{0.13,0.29,0.53}{#1}}
\newcommand{\DecValTok}[1]{\textcolor[rgb]{0.00,0.00,0.81}{#1}}
\newcommand{\DocumentationTok}[1]{\textcolor[rgb]{0.56,0.35,0.01}{\textbf{\textit{#1}}}}
\newcommand{\ErrorTok}[1]{\textcolor[rgb]{0.64,0.00,0.00}{\textbf{#1}}}
\newcommand{\ExtensionTok}[1]{#1}
\newcommand{\FloatTok}[1]{\textcolor[rgb]{0.00,0.00,0.81}{#1}}
\newcommand{\FunctionTok}[1]{\textcolor[rgb]{0.13,0.29,0.53}{\textbf{#1}}}
\newcommand{\ImportTok}[1]{#1}
\newcommand{\InformationTok}[1]{\textcolor[rgb]{0.56,0.35,0.01}{\textbf{\textit{#1}}}}
\newcommand{\KeywordTok}[1]{\textcolor[rgb]{0.13,0.29,0.53}{\textbf{#1}}}
\newcommand{\NormalTok}[1]{#1}
\newcommand{\OperatorTok}[1]{\textcolor[rgb]{0.81,0.36,0.00}{\textbf{#1}}}
\newcommand{\OtherTok}[1]{\textcolor[rgb]{0.56,0.35,0.01}{#1}}
\newcommand{\PreprocessorTok}[1]{\textcolor[rgb]{0.56,0.35,0.01}{\textit{#1}}}
\newcommand{\RegionMarkerTok}[1]{#1}
\newcommand{\SpecialCharTok}[1]{\textcolor[rgb]{0.81,0.36,0.00}{\textbf{#1}}}
\newcommand{\SpecialStringTok}[1]{\textcolor[rgb]{0.31,0.60,0.02}{#1}}
\newcommand{\StringTok}[1]{\textcolor[rgb]{0.31,0.60,0.02}{#1}}
\newcommand{\VariableTok}[1]{\textcolor[rgb]{0.00,0.00,0.00}{#1}}
\newcommand{\VerbatimStringTok}[1]{\textcolor[rgb]{0.31,0.60,0.02}{#1}}
\newcommand{\WarningTok}[1]{\textcolor[rgb]{0.56,0.35,0.01}{\textbf{\textit{#1}}}}
\usepackage{graphicx}
\makeatletter
\newsavebox\pandoc@box
\newcommand*\pandocbounded[1]{% scales image to fit in text height/width
  \sbox\pandoc@box{#1}%
  \Gscale@div\@tempa{\textheight}{\dimexpr\ht\pandoc@box+\dp\pandoc@box\relax}%
  \Gscale@div\@tempb{\linewidth}{\wd\pandoc@box}%
  \ifdim\@tempb\p@<\@tempa\p@\let\@tempa\@tempb\fi% select the smaller of both
  \ifdim\@tempa\p@<\p@\scalebox{\@tempa}{\usebox\pandoc@box}%
  \else\usebox{\pandoc@box}%
  \fi%
}
% Set default figure placement to htbp
\def\fps@figure{htbp}
\makeatother
\setlength{\emergencystretch}{3em} % prevent overfull lines
\providecommand{\tightlist}{%
  \setlength{\itemsep}{0pt}\setlength{\parskip}{0pt}}
\usepackage{bookmark}
\IfFileExists{xurl.sty}{\usepackage{xurl}}{} % add URL line breaks if available
\urlstyle{same}
\hypersetup{
  pdftitle={Prior Specification},
  hidelinks,
  pdfcreator={LaTeX via pandoc}}

\title{Prior Specification}
\author{}
\date{\vspace{-2.5em}2026-04-17}

\begin{document}
\maketitle

\subsection{Introduction}\label{introduction}

Choosing a prior is the most consequential modelling decision in any
Bayesian analysis after the choice of likelihood. The prior shapes the
posterior strongly when data are sparse, almost imperceptibly when data
are abundant, and unpredictably in the middle ground that most applied
work occupies. Three categories cover most needs: flat or improper
priors, weakly informative priors that regularise without imposing
substantive content, and informative priors that incorporate genuine
pre-data knowledge. Modern Bayesian practice has converged on weakly
informative priors as the default because they stabilise sampling,
regularise small-sample inference, and remain agnostic enough that the
data dominate the posterior whenever they are informative.

\subsection{Prerequisites}\label{prerequisites}

A working understanding of Bayes' theorem, the role of the prior in
posterior updating, and basic familiarity with parameter scales
(standardised vs.~raw, log-scale vs.~linear).

\subsection{Theory}\label{theory}

\textbf{Flat / improper priors} assign constant density across the
support. They sometimes appear intuitive but can yield improper
posteriors in hierarchical models or near boundary points, and they are
not invariant under reparameterisation.

\textbf{Weakly informative priors} are proper distributions broad enough
to admit any plausible value while ruling out absurd ones. Standard
choices are \(\mathcal N(0, 2.5)\) for standardised slopes,
\(\mathcal N(0, 10)\) for intercepts,
\(\mathrm{half\text{-}Cauchy}(0, 5)\) or
\(\mathrm{half\text{-}Normal}(0, 2.5)\) for variance components, and
\(\mathrm{LKJ}(2)\) for correlation matrices. They regularise estimates
without injecting domain claims.

\textbf{Informative priors} carry substantive content from previous
studies, theoretical bounds, or expert elicitation. They are appropriate
when the information is defensible and reproducible, and inappropriate
when chosen post-hoc to push the posterior somewhere convenient.

The \textbf{prior predictive check} --- simulating from the prior,
propagating through the likelihood, and inspecting the implied range of
\texttt{y} --- is the single most useful sanity check. If the prior
implies blood pressures of 1,000 mm Hg or odds ratios of \(10^6\), it is
too wide.

\subsection{Assumptions}\label{assumptions}

The prior reflects pre-data information honestly, the likelihood is
correctly specified, and the prior is proper unless the resulting
posterior has been verified to be proper analytically.

\subsection{R Implementation}\label{r-implementation}

\begin{Shaded}
\begin{Highlighting}[]
\FunctionTok{library}\NormalTok{(brms)}

\CommentTok{\# Logistic regression with weakly informative priors}
\NormalTok{priors }\OtherTok{\textless{}{-}} \FunctionTok{prior}\NormalTok{(}\FunctionTok{normal}\NormalTok{(}\DecValTok{0}\NormalTok{, }\FloatTok{2.5}\NormalTok{), }\AttributeTok{class =} \StringTok{"b"}\NormalTok{) }\SpecialCharTok{+}
          \FunctionTok{prior}\NormalTok{(}\FunctionTok{normal}\NormalTok{(}\DecValTok{0}\NormalTok{, }\DecValTok{10}\NormalTok{),  }\AttributeTok{class =} \StringTok{"Intercept"}\NormalTok{)}

\CommentTok{\# brms: inspect priors before fitting}
\FunctionTok{get\_prior}\NormalTok{(y }\SpecialCharTok{\textasciitilde{}}\NormalTok{ x1 }\SpecialCharTok{+}\NormalTok{ x2, }\AttributeTok{data =} \FunctionTok{data.frame}\NormalTok{(}\AttributeTok{y =} \FunctionTok{rbinom}\NormalTok{(}\DecValTok{100}\NormalTok{, }\DecValTok{1}\NormalTok{, }\FloatTok{0.5}\NormalTok{),}
                                          \AttributeTok{x1 =} \FunctionTok{rnorm}\NormalTok{(}\DecValTok{100}\NormalTok{), }\AttributeTok{x2 =} \FunctionTok{rnorm}\NormalTok{(}\DecValTok{100}\NormalTok{)),}
          \AttributeTok{family =}\NormalTok{ bernoulli)}
\end{Highlighting}
\end{Shaded}

\subsection{Output \& Results}\label{output-results}

\texttt{get\_prior()} lists every parameter class with its default
prior, ready to be overridden. After fitting,
\texttt{prior\_summary(fit)} confirms exactly which priors entered the
model --- a step worth running for any analysis intended for
publication.

\subsection{Interpretation}\label{interpretation}

A reporting sentence: ``We placed weakly informative priors on all
coefficients (\(\beta \sim \mathcal N(0, 2.5)\) on standardised
predictors, intercept \(\sim \mathcal N(0, 10)\)) and a
\(\mathrm{half\text{-}Cauchy}(0, 2.5)\) prior on the random-effect
standard deviation; a sensitivity analysis under tighter
\(\mathcal N(0, 1)\) slopes produced indistinguishable posterior
summaries.'' Always disclose priors and at least one sensitivity check;
reviewers increasingly expect both.

\subsection{Practical Tips}\label{practical-tips}

\begin{itemize}
\tightlist
\item
  Defaults from \texttt{brms} and \texttt{rstanarm} are weakly
  informative and well-engineered; override them only when you have a
  reason.
\item
  Flat priors are not ``neutral''; they can dominate hierarchical
  posteriors at the boundary and make sampling unstable.
\item
  For variance components, use half-Cauchy or half-Normal priors with
  scales matched to the natural variability of the data; never use a
  flat prior on a SD.
\item
  For unstandardised predictors, multiply prior scales by the empirical
  SD of the predictor so that the implied effect range is plausible.
\item
  Run prior predictive simulations (\texttt{sample\_prior\ =\ "only"} in
  \texttt{brms}) before fitting; the implied range of \texttt{y} is the
  simplest check that the prior is not absurd.
\item
  Sensitivity analysis is mandatory for small-sample studies; report the
  posterior under at least two prior scales differing by an order of
  magnitude.
\end{itemize}

\subsection{R Packages Used}\label{r-packages-used}

\texttt{brms} for prior specification with \texttt{get\_prior()} /
\texttt{prior\_summary()}, \texttt{rstanarm} for an alternative
interface with auto-scaled defaults, \texttt{bayesplot} for prior
predictive visualisation, and \texttt{bayestestR} for tidy summaries
that report priors alongside posteriors.

\subsection{For Reviewers}\label{for-reviewers}

\textbf{What to look for in a paper using this method.}

\begin{itemize}
\tightlist
\item
  \textbf{Common misapplications.}
\item
  \textbf{Diagnostics that should be reported but often aren't.}
\item
  \textbf{Red flags in tables and figures.}
\item
  \textbf{What to verify.}
\item
  \textbf{What an adequate Methods paragraph must contain.}
\end{itemize}

\subsection{See also --- labs in this
chapter}\label{see-also-labs-in-this-chapter}

\begin{itemize}
\tightlist
\item
  \href{labs/lab-bayesian-thinking-control-vs-decision-error.Rmd}{Bayesian
  thinking; α control vs decision error}
\item
  \href{labs/lab-brms-stan-loo-hierarchical-models.Rmd}{brms / Stan,
  LOO, hierarchical models}
\end{itemize}

\end{document}
